% GENERATED FILE. Edit canonical lessons, then run npm run generate:books.
% canonical-source-hash: fnv1a64:eb992e8a481018fd
% canonical-lessons: RU-C422-vylet, RU-C422-vyletat, RU-C422-prilyot, RU-C422-terminal, RU-C422-styuardessa

\chapter{At the airport}
\label{ch:ru-at-the-airport}
\hlchaptermodality{422}

\begin{chapteropening}
\textbf{By the end of this chapter:} \emph{I can read the airport board: departure, flying out, arrival, the terminal and the flight attendant.}

Complete the last lesson of chapter 422: I can read the airport board: departure, flying out, arrival, the terminal and the flight attendant.
\end{chapteropening}

\section[\texorpdfstring{\ru{вылет} (výlet)}{вылет (výlet)}]{\ru{вылет} (výlet) — a departure (of a flight)}

\label{lesson:RU-C422-vylet}



\begin{quote}
Before the new one: say the Russian for a porthole, then the Russian for the steps.
\end{quote}

\subsection*{You'll want to know: \ru{вылет}}
\textbf{\ru{вылет}} (\emph{výlet}) — \textquotedblleft{}a departure (of a flight)\textquotedblright{}.

\begin{quote}
\ru{Вылет} \ru{задержали} \ru{на} \ru{час}.
\end{quote}

\begin{quote}
\emph{Výlet zaderzháli na chas.}
\end{quote}

\begin{quote}
The departure was delayed by an hour.
\end{quote}

Masculine; the stress is on \emph{\ru{вы}-} (\emph{vý-}). \emph{\ru{Время} \ru{вылета}} — departure time.

\begin{cousinweb}
\emph{\ru{вылет}} is built from \emph{\ru{вы}-} 'out' and \emph{\ru{лёт}} 'flight', the root of \emph{\ru{летать}} 'to fly'. EtymDB lists \emph{\ru{взлёт}} from the same root, Proto-Slavic \emph{vŭzletŭ}.
\end{cousinweb}

\begin{grammarlens}[title={What you've built}]
One more word for telling, explaining and describing at B1.
\end{grammarlens}

\subsection*{Your turn}
\begin{itemize}
\raggedright
  \item \emph{Say it:} \emph{výlet}
  \item \emph{Say it:} \emph{výlet}, once more
  \item \emph{Say it:} a sentence of your own with \emph{výlet}
  \item \emph{Recall:} say the Russian for a porthole, then the Russian for the steps, then say \emph{výlet} again
\end{itemize}

\begin{tcolorbox}[breakable,colback=teal!4,colframe=teal!35!black,title={Before you move on}]
What does \ru{вылет} mean? (\textquotedblleft{}A departure (of a flight)\textquotedblright{}.) Say it once more.
\end{tcolorbox}
\section[\texorpdfstring{\ru{вылетать} (vyletát)}{вылетать (vyletát)}]{\ru{вылетать} (vyletát) — to fly out, to depart by plane}

\label{lesson:RU-C422-vyletat}



\begin{quote}
Before the new one: say the Russian for the steps, then the Russian for a departure.
\end{quote}

\subsection*{You'll want to know: \ru{вылетать}}
\textbf{\ru{вылетать}} (\emph{vyletát}) — \textquotedblleft{}to fly out, to depart by plane\textquotedblright{}.

\begin{quote}
\ru{Мы} \ru{вылетаем} \ru{рано} \ru{утром}.
\end{quote}

\begin{quote}
\emph{My vyletáyem ráno útrom.}
\end{quote}

\begin{quote}
We fly out early in the morning.
\end{quote}

Imperfective; perfective \emph{\ru{вылететь}} (\emph{\ru{вылечу}, \ru{вылетишь}}). \emph{\ru{Взлетать}} is the take-off itself; \emph{\ru{вылетать}} is leaving by air.

\begin{cousinweb}
\emph{\ru{вылетать}} is built from \emph{\ru{вы}-} 'out' and \emph{\ru{летать}} 'to fly'. EtymDB lists \emph{\ru{взлёт}} from the same root, Proto-Slavic \emph{vŭzletŭ}.
\end{cousinweb}

\begin{grammarlens}[title={What you've built}]
One more word for telling, explaining and describing at B1.
\end{grammarlens}

\subsection*{Your turn}
\begin{itemize}
\raggedright
  \item \emph{Say it:} \emph{vyletát}
  \item \emph{Say it:} \emph{vyletát}, once more
  \item \emph{Say it:} a sentence of your own with \emph{vyletát}
  \item \emph{Recall:} say the Russian for the steps, then the Russian for a departure, then say \emph{vyletát} again
\end{itemize}

\begin{tcolorbox}[breakable,colback=teal!4,colframe=teal!35!black,title={Before you move on}]
What does \ru{вылетать} mean? (\textquotedblleft{}To fly out, to depart by plane\textquotedblright{}.) Say it once more.
\end{tcolorbox}
\section[\texorpdfstring{\ru{прилёт} (prilyót)}{прилёт (prilyót)}]{\ru{прилёт} (prilyót) — an arrival (of a flight)}

\label{lesson:RU-C422-prilyot}



\begin{quote}
Before the new one: say the Russian for a departure, then the Russian for to fly out.
\end{quote}

\subsection*{You'll want to know: \ru{прилёт}}
\textbf{\ru{прилёт}} (\emph{prilyót}) — \textquotedblleft{}an arrival (of a flight)\textquotedblright{}.

\begin{quote}
\ru{После} \ru{прилёта} \ru{мы} \ru{сразу} \ru{поехали} \ru{в} \ru{гостиницу}.
\end{quote}

\begin{quote}
\emph{Pósle prilyóta my srázu poyékhali v gostínitsu.}
\end{quote}

\begin{quote}
After we landed we went straight to the hotel.
\end{quote}

Masculine. \emph{\ru{Зал} \ru{прилёта}} — the arrivals hall.

\begin{cousinweb}
\emph{\ru{прилёт}} is built from \emph{\ru{при}-} 'to, towards' and \emph{\ru{лёт}} 'flight', the root of \emph{\ru{летать}} 'to fly'. EtymDB lists \emph{\ru{взлёт}} from the same root, Proto-Slavic \emph{vŭzletŭ}.
\end{cousinweb}

\begin{grammarlens}[title={What you've built}]
One more word for telling, explaining and describing at B1.
\end{grammarlens}

\subsection*{Your turn}
\begin{itemize}
\raggedright
  \item \emph{Say it:} \emph{prilyót}
  \item \emph{Say it:} \emph{prilyót}, once more
  \item \emph{Say it:} a sentence of your own with \emph{prilyót}
  \item \emph{Recall:} say the Russian for a departure, then the Russian for to fly out, then say \emph{prilyót} again
\end{itemize}

\begin{tcolorbox}[breakable,colback=teal!4,colframe=teal!35!black,title={Before you move on}]
What does \ru{прилёт} mean? (\textquotedblleft{}An arrival (of a flight)\textquotedblright{}.) Say it once more.
\end{tcolorbox}
\section[\texorpdfstring{\ru{терминал} (terminál)}{терминал (terminál)}]{\ru{терминал} (terminál) — a terminal (at an airport); a card machine}

\label{lesson:RU-C422-terminal}



\begin{quote}
Before the new one: say the Russian for to fly out, then the Russian for an arrival.
\end{quote}

\subsection*{You'll want to know: \ru{терминал}}
\textbf{\ru{терминал}} (\emph{terminál}) — \textquotedblleft{}a terminal (at an airport); a card machine\textquotedblright{}.

\begin{quote}
\ru{Между} \ru{терминалами} \ru{ходит} \ru{автобус}.
\end{quote}

\begin{quote}
\emph{Mézhdu terminálami khódit avtóbus.}
\end{quote}

\begin{quote}
A bus runs between the terminals.
\end{quote}

Masculine. Also a card reader in a shop: \emph{\ru{оплатить} \ru{через} \ru{терминал}}.

\begin{cousinweb}
\emph{\ru{терминал}} is usually traced to English \emph{terminal}, which EtymDB derives through French from Late Latin \emph{terminālis}, on \emph{terminus} 'a boundary, an end'. The Russian word is not in EtymDB.
\end{cousinweb}

\begin{grammarlens}[title={What you've built}]
One more word for telling, explaining and describing at B1.
\end{grammarlens}

\subsection*{Your turn}
\begin{itemize}
\raggedright
  \item \emph{Say it:} \emph{terminál}
  \item \emph{Say it:} \emph{terminál}, once more
  \item \emph{Say it:} a sentence of your own with \emph{terminál}
  \item \emph{Recall:} say the Russian for to fly out, then the Russian for an arrival, then say \emph{terminál} again
\end{itemize}

\begin{tcolorbox}[breakable,colback=teal!4,colframe=teal!35!black,title={Before you move on}]
What does \ru{терминал} mean? (\textquotedblleft{}A terminal (at an airport); a card machine\textquotedblright{}.) Say it once more.
\end{tcolorbox}
\section[\texorpdfstring{\ru{стюардесса} (styuardéssa)}{стюардесса (styuardéssa)}]{\ru{стюардесса} (styuardéssa) — a flight attendant (woman), a stewardess}

\label{lesson:RU-C422-styuardessa}



\begin{quote}
Before the new one: say the Russian for an arrival, then the Russian for a terminal.
\end{quote}

\subsection*{You'll want to know: \ru{стюардесса}}
\textbf{\ru{стюардесса}} (\emph{styuardéssa}) — \textquotedblleft{}a flight attendant (woman), a stewardess\textquotedblright{}.

\begin{quote}
\ru{Стюардесса} \ru{принесла} \ru{нам} \ru{воду}.
\end{quote}

\begin{quote}
\emph{Styuardéssa prineslá nam vódu.}
\end{quote}

\begin{quote}
The flight attendant brought us water.
\end{quote}

The official word is \emph{\ru{бортпроводница}}; a man is a \emph{\ru{стюард}}.

\begin{cousinweb}
\emph{\ru{стюардесса}} was borrowed from English \emph{stewardess} (EtymDB).
\end{cousinweb}

\begin{grammarlens}[title={What you've built}]
One more word for telling, explaining and describing at B1.
\end{grammarlens}

\subsection*{Your turn}
\begin{itemize}
\raggedright
  \item \emph{Say it:} \emph{styuardéssa}
  \item \emph{Say it:} \emph{styuardéssa}, once more
  \item \emph{Say it:} a sentence of your own with \emph{styuardéssa}
  \item \emph{Recall:} say the Russian for an arrival, then the Russian for a terminal, then say \emph{styuardéssa} again
\end{itemize}

\begin{tcolorbox}[breakable,colback=teal!4,colframe=teal!35!black,title={Before you move on}]
What does \ru{стюардесса} mean? (\textquotedblleft{}A flight attendant (woman), a stewardess\textquotedblright{}.) Say it once more.
\end{tcolorbox}
